\documentclass{ximera}
\title{Velocity}

\newcommand{\pskip}{\vskip 0.1 in}

\begin{document}
\begin{abstract}
The velocity vector.
\end{abstract}
\maketitle



\section{The Velocity Vector}
\emph{Velocity} is the rate of change of position with respect to time. Since position is relative, so is velocity. 

There are two types of velocity. Average velocity and instantaneous vecocity. You calculate average velocity as you would any average rate of change, here as the change in position divided by the change in time.

Calculating instantaneous velocity usually requires ideas of differential calculus. But in the special case of uniform circular motion, we do not need calculus.

The key ideas are these. At any moment, the instantaneous velocity of a point mass points in the direction of motion and has length equal to the speed of the mass.

For a beetle moving around a circle at a constant speed, because the velocity vector is tangent to the path, it is always perpendicular to the position vector of the beetle relative to the center of its path.

\begin{question} \label{Qe3r3rbnnn}
\begin{enumerate}
\item A beetle crawls counterclockwise around a circle centered at the point $A(3,1)$ at a constant speed of $12$ cm/sec. Find the velocity of the beetle as it passes the point $P(-8,4)$. All coordinates are measured in cm.

\item As an ant rotates about the origin at a constant rate of $1/5$ rad/sec as it crawls clockwise around a circle centered at the origin. Find the velocity of the ant as it passes the point $(7,-10)$.  All coordinates are measured in cm.

\end{enumerate}

\begin{explanation}
\begin{enumerate}
\item Since the velocity vector is tangent to the circular path and in the direction of motion, we can find its direction by rotating the position vector (coordinates in cm)
\begin{align*}
  \overrightarrow{AP} &= \overrightarrow{\answer{OP}} - \overrightarrow{\answer{OA}} \\
                                &= \langle \answer{-8}, \answer{4} \rangle - \langle \answer{3} , \answer{1} \rangle  \\
                                &= \langle  \answer{-11} , \answer{3} \rangle  
\end{align*}
$\pi/2$ radians counterclockwise. This gives the vector
\[
 {\bf w} = \langle \answer{-3} , \answer{-11} \rangle \text{ cm}.
\]
Since the velocity vector ${\bf v}$ points in the same direction as ${\bf w}$ and has length $v = 2$ cm/sec equal to the speed, we just need to scale ${\bf w}$ by the appropriate factor to get ${\bf v}$. This factor is
\begin{align*}
    \lambda  &= \frac{v}{|{\bf w}|}  \\
                  &= \frac{\answer{12} \frac{\text{cm}}{\text{sec}}}{\answer{\sqrt{130}\text{ cm}}} \\
                 &= \answer{\frac{12}{\sqrt{130}}} \text{ sec}^{-1} .
\end{align*}

So the beetle's velocity as it passes the point $P$ is
\begin{align*}
   {\bf v}  &= {\bf w} \left( \frac{v}{|{\bf w}|} \right)  \\
               &= \left \langle  \answer{\frac{-36}{\sqrt{130}}} ,  \answer{\frac{-132}{\sqrt{130}}}\right \rangle \frac{\text{cm}}{\text{sec}}.
\end{align*}

\end{enumerate}
\end{explanation}
\end{question}

\end{document}
